\subsection{The interface with the technical supplement}
\label{sec:finite-interface}
The exact inverse in Section~\ref{sec:two-field-rigidity} is proved
entirely in this article. For the finite construction we use the
auxiliary estimates in Table~\ref{tab:finite-interface}. We state their
inputs and outputs here to specify the dependency of the argument.
All constants in this subsection depend only on the fixed priors.
The supplement proves the indicated geometric or analytic estimates;
the observation procedures in the present article supply their inputs.

\begin{table}[ht]
\centering
\small
\caption{Auxiliary estimates and their proof locations in the supplement.}
\label{tab:finite-interface}
\begin{tabular}{@{}p{.48\textwidth}p{.46\textwidth}@{}}
\toprule
Estimate used here & Proof in the supplement\\
\midrule
Overlap mass and a boundary layer &
Lemma~\ref{H-lem:cap-mass}\\[3pt]
Complete coarse components and radial brackets &
Lemma~\ref{H-lem:coarse}\\[3pt]
Fixed-accuracy projection normals &
Lemma~\ref{H-lem:rare-coarse-normals}\\[3pt]
Bisection and radial interpolation &
Lemmas~\ref{H-lem:bisection} and \ref{H-lem:radial}\\[3pt]
Signed support smoothing &
Section~\ref{H-sec:stationary}, proof of
Proposition~\ref{H-prop:jitter-modulus}\\[3pt]
Period decisions and common-translation invariance &
Lemmas~\ref{H-lem:recognition}, \ref{H-lem:locking}, and
\ref{H-lem:registration-period-margin}\\
\bottomrule
\end{tabular}
\end{table}

\paragraph{Mass, coarse components and normals.}
Let $P=C+(-A)$, with the smoothness, curvature and density bounds
\eqref{eq:two-field-geometric-priors} and
\eqref{eq:two-field-finite-priors}. There are known $c,e_0>0$ such that
\[
 x\in P_{-e}\ \Longrightarrow\
 \int\one_C(x+z)j(z)\dd z\ge c e^{\gamma+3/2},
 \qquad 0<e<e_0.
\]
The constants are uniform when $A$ varies over its stated class.
For a family of separated convex bodies of bounded diameter and
positive rolling radius, a fixed sufficiently fine grid of labels
which retain every grid node in an inner parallel body and discard every
grid node outside the body gives one approximate hull per complete
component. The construction supplies a center $o$ and fixed
$0<r<R$ with
$\overline B_r(o)\subset P\cap\widehat P\subset\overline B_R(o)$,
and isolated radial intervals enclosing every boundary point.
The aperture has a collar larger than twice the component diameter,
with additional clustering and numerical reserves; discarding clusters meeting the outer cutoff then
keeps every required complete component and no fragment. Refining the
radial data to any prescribed fixed accuracy supplies projection
normals to that accuracy throughout fixed tubular neighborhoods of
the intervals. The interpolation proof permits any fixed normal
tolerance; its displayed value $1/32$ is one choice.
In this article the mass bound and the prefix inverse furnish these
labels, with the threshold and its error reserve specified below.

\paragraph{Bisection and interpolation.}
Suppose a radial interval has width $w_0$ and a label can differ from
the true inside/outside label only within distance $Ae$ of the true
radial endpoint. The label may depend on earlier tests and need not
be monotone. After $k$ midpoint tests the returned radial error is
at most $Ae+2^{-k-1}w_0$. For a body with the common inball,
curvature and $C^{6,\beta}$ support bounds, its radial function has
a uniform $C^{6,\beta}$ bound. From angular spacing $h$ and value
errors at most $e$, the supplement constructs a smooth interpolation
with
\[
 \|\widehat R-R\|_{C^j}\le C(h^{s-j}+eh^{-j}),
 \qquad 0\le j\le3.
\]
When $e\le C_0h^s$, the resulting curve is strictly convex for small
$h$; its Hausdorff error is $Ch^s$ and its support $C^2$ error is
$Ch^{s-2}$. These assertions depend on the stated label and geometric
inputs, irrespective of how the labels were acquired.

\paragraph{Support smoothing.}
There is a fixed signed smooth kernel $\kappa$ on a bounded interval
with integral one and vanishing moments of orders one through six.
For its periodized dilation $T_b$, a function with uniformly bounded
$C^{6,\beta}$ norm satisfies
\[
 \|T_bf-f\|_{C^j}\le Cb^{s-j},\qquad
 \|T_bg\|_{C^j}\le Cb^{-j}\|g\|_\infty,
 \qquad 0\le j\le3.
\]
The supplement gives the explicit dilation coefficients
$8/5,-4/5,8/35,-1/35$ at scales $1,2,3,4$. Taylor's formula and
the moment cancellations give the first estimate, and the $L^1$
norms of the differentiated kernel give the second. The kernel is
used for numerical integration of support functions.

\paragraph{Period locking.}
For the bounded presentation \eqref{eq:two-field-presentation},
put $B=1+L_0$ and $V_*=V_0/r_0$. Complete component estimates with
matched Hausdorff error $e$ in the protected square $Q_{10B}$ suffice
for the following decisions. Under the patch margin
\eqref{eq:two-field-patch-margin}, and below a known positive
threshold for $e$, the tests accept exactly those candidates whose
matched true differences are periods. The accepted true
vectors generate $\Pi$, and $\covol\Pi\ge V_*$. If their lengths
are bounded by the prior-dependent $U_0$, coordinates in an accepted
independent-pair basis have reduced denominators at most
$Q=\max(1,\lceil U_0^2/V_*\rceil)$. Comparison below the rational
gap $Q^{-2}$ recovers the exact relations, the primitive orbit
count, a matched real period basis with error $Ce$, and primitive
free area with error $Ce$. The support comparison and the orbitwise
defect are unchanged by adding one common Minkowski summand or by
one common translation, provided the expanded components remain
separated. Thus the same $\eta$ applies in the canonical frame.
The numerical enclosures, cutoff reserves and periodic assembly
are included in the cited proofs. The finite-configuration result
uses the given complete aperture and omits this period step.

The two-command construction below therefore needs neither the
directional-germ data nor the reciprocal randomized preparations of
the other experiments in the supplement. Each auxiliary hypothesis
has been stated here, and each observation required for the current
finite theorem is supplied by its own two fixed forward channels.
